\documentclass[12pt]{article}
\usepackage[utf8]{inputenc}
\usepackage[T1]{fontenc}
\usepackage[margin=1in]{geometry}
\usepackage{setspace}
\usepackage{xcolor}
\usepackage[colorlinks=true, urlcolor=blue!60!black, linkcolor=blue!60!black]{hyperref}
\setlength{\parindent}{0em}
\setlength{\parskip}{0.6em}
\renewcommand{\baselinestretch}{1.0}
\begin{document}
\begin{center}
{\LARGE \textit{Friends of the Earth, Inc.\ v.\ Laidlaw Environmental Services (TOC), Inc.}, 528 U.S.\ 167 (2000)\par}
\vspace{0.5em}
{\large Case Brief --- Two-Page Edition\par}
\end{center}

\section{Memory Jogger}
Citizen-suit plaintiffs who curtailed use of a polluted river had standing; deterrent penalties redress ongoing violations; and voluntary cessation rarely moots a case (\href{https://supreme.justia.com/cases/federal/us/528/167/}{Justia}).

\section{Detailed Case Facts}
Laidlaw's South Carolina incinerator committed nearly 500 violations of its mercury discharge permit. Environmental groups sued under the Clean Water Act for civil penalties payable to the Treasury after Laidlaw arranged a quick, lenient state settlement.

\section{Procedural History}
The district court found no proof of environmental harm, denied an injunction because Laidlaw had come into compliance, and imposed a \$405,800 penalty. The Fourth Circuit held the case moot under \textit{Steel Co.} Laidlaw then closed the plant but kept its permit. The Supreme Court reversed.

\section{Judicial Votes}
Reversed 7--2. Ginsburg wrote for Rehnquist, Stevens, O'Connor, Kennedy, Souter, and Breyer; Stevens and Kennedy concurred; Scalia (with Thomas) dissented.

\section{Holding}
Members' reasonable concerns that curtailed their use of the river are injury in fact; penalties for violations ongoing at filing redress that injury through deterrence; and a defendant claiming voluntary cessation must show that recurrence cannot reasonably be expected.

\section{Analysis of Opinions}
\textbf{Majority:} the injury that matters is to the plaintiff, not the environment; deterrence redresses ongoing violations, unlike the wholly past violations in \textit{Steel Co.}; and mootness differs from standing in burden and purpose. \textbf{Kennedy:} private suits for public fines raise Article~II questions. \textbf{Scalia:} injury became a ``mere pleading requirement,'' and Treasury penalties offer no personal redress.

\section{Examples: Future Applications}
\textbf{Same-Side:} a plant that complies mid-suit cannot easily moot the case. \textbf{Opposite-Side:} violations that ended before the suit was filed. \textbf{Fence-Sitter:} a closed plant that keeps its permit.

\section{Critique}
Critics say standing now rests on subjective fear of technical violations, and that citizen suits for public fines intrude on executive enforcement. \textit{Laidlaw} also shows that mootness is not simply standing in a time frame.

\section{Key Quotations}
\begin{quote}
``[A] defendant claiming that its voluntary compliance moots a case bears the formidable burden of showing that it is absolutely clear the allegedly wrongful behavior could not reasonably be expected to recur.''
\end{quote}
\noindent --- Justice Ginsburg, for the Court, 528 U.S.\ at 190.

\end{document}
